\chapter{lambda、函数对象和可调用设计}

\section{问题入口：排序规则从哪里来}

给用户列表排序时，排序规则常常不是固定的。按 ID 升序、按姓名升序、按活跃状态优先，都是合理需求。最直接的写法是为每个规则写一个函数：

\begin{lstlisting}[language=C++,caption={普通比较函数}]
bool by_id(const User& lhs, const User& rhs)
{
    return lhs.id < rhs.id;
}

std::ranges::sort(users, by_id);
\end{lstlisting}

这个版本没问题。但如果排序规则只在一处使用，单独起一个函数名可能让代码跳来跳去。lambda 提供了就地定义小函数的方式：

\begin{lstlisting}[language=C++,caption={lambda 就地表达比较规则}]
std::ranges::sort(users, [](const User& lhs, const User& rhs) {
    return lhs.id < rhs.id;
});
\end{lstlisting}

lambda 的价值不是“写得短”，而是让局部规则贴近使用位置。读者看到排序调用时，就能看到排序规则。

\section{捕获：lambda 如何拿到外部变量}

假设要筛选分数大于阈值的学生。阈值来自外部变量：

\begin{lstlisting}[language=C++,caption={按值捕获阈值}]
int threshold = 60;
auto passed = [threshold](const Student& student) {
    return student.score >= threshold;
};
\end{lstlisting}

\code{[threshold]} 表示把当前 \code{threshold} 的值复制进 lambda 对象。之后外部变量变化，不影响这个 lambda 已经保存的值。按引用捕获则不同：

\begin{lstlisting}[language=C++,caption={按引用捕获阈值}]
int threshold = 60;
auto passed = [&threshold](const Student& student) {
    return student.score >= threshold;
};
threshold = 80;
\end{lstlisting}

这时 lambda 每次使用的是外部 \code{threshold} 当前值。按引用捕获更灵活，也更危险：如果 lambda 活得比被捕获变量更久，就会悬空。

\section{捕获生命周期反例}

下面的函数返回一个 lambda，但它引用了局部变量：

\begin{lstlisting}[language=C++,caption={返回引用局部变量的 lambda 是错误}]
auto make_bad_filter()
{
    int threshold = 60;
    return [&threshold](const Student& student) {
        return student.score >= threshold;
    };
}
\end{lstlisting}

函数返回后，\code{threshold} 已销毁。返回的 lambda 保存的是悬空引用。正确做法是按值捕获：

\begin{lstlisting}[language=C++,caption={返回 lambda 时按值保存状态}]
auto make_score_filter(int threshold)
{
    return [threshold](const Student& student) {
        return student.score >= threshold;
    };
}
\end{lstlisting}

这个反例推出捕获规则：lambda 不逃出当前作用域时，引用捕获可以很方便；lambda 要保存、返回、放进异步任务或跨线程使用时，优先按值捕获必要状态。

\section{函数对象：有名字的可调用类型}

lambda 本质上会生成一个匿名函数对象。你也可以自己写函数对象：

\begin{lstlisting}[language=C++,caption={带状态的函数对象}]
class ScoreAtLeast {
public:
    explicit ScoreAtLeast(int threshold)
        : threshold_{threshold}
    {
    }

    [[nodiscard]] bool operator()(const Student& student) const noexcept
    {
        return student.score >= this->threshold_;
    }

private:
    int threshold_ = 0;
};
\end{lstlisting}

使用方式和 lambda 一样：

\begin{lstlisting}[language=C++,caption={函数对象作为谓词}]
const auto passed = ScoreAtLeast{60};
const bool ok = passed(student);
\end{lstlisting}

什么时候写函数对象？当规则需要复用、需要命名、需要多个成员函数、或者测试这个规则本身时，具名类型比一个很长的 lambda 更清楚。

\section{std::function 的边界}

\code{std::function} 可以保存任意满足签名的可调用对象：

\begin{lstlisting}[language=C++,caption={用 std::function 保存回调}]
using Filter = std::function<bool(const Student&)>;

std::vector<Student> filter_students(std::span<const Student> students, const Filter& filter)
{
    std::vector<Student> result;
    for (const Student& student : students) {
        if (filter(student)) {
            result.push_back(student);
        }
    }
    return result;
}
\end{lstlisting}

\code{std::function} 的好处是接口稳定，可以在运行期保存不同类型的 callable。代价是类型擦除可能带来间接调用和分配。热路径上，如果规则类型在编译期已知，可以用模板：

\begin{lstlisting}[language=C++,caption={模板接受任意谓词}]
template <typename Filter>
std::vector<Student> filter_students_fast(std::span<const Student> students, Filter filter)
{
    std::vector<Student> result;
    for (const Student& student : students) {
        if (filter(student)) {
            result.push_back(student);
        }
    }
    return result;
}
\end{lstlisting}

两者没有绝对优劣。需要运行期多态和稳定 ABI 时，\code{std::function} 合理；需要极低开销和内联机会时，模板更合适。

\section{回调接口不要制造隐藏所有权}

如果回调捕获了对象指针，就要考虑对象生命周期：

\begin{lstlisting}[language=C++,caption={回调捕获裸指针的风险}]
class Button {
public:
    void on_click(std::function<void()> callback)
    {
        this->callback_ = std::move(callback);
    }

private:
    std::function<void()> callback_;
};
\end{lstlisting}

如果调用者写：

\begin{lstlisting}[language=C++,caption={可能悬空的 this 捕获}]
button.on_click([this] {
    this->save();
});
\end{lstlisting}

那么必须保证回调执行时当前对象还活着。GUI、异步任务、线程池中，这个保证经常很难。可以考虑捕获 \code{std::weak_ptr}，执行前尝试锁定；或者让拥有者负责在析构前注销回调。回调设计本质上是生命周期设计。

\section{完整案例：可配置验证流水线}

假设要验证注册用户：用户名合法、年龄达到要求、邮箱域名允许。每个规则都是一个谓词：

\begin{lstlisting}[language=C++,caption={验证规则流水线}]
struct Registration {
    std::string username;
    int age = 0;
    std::string email;
};

using Rule = std::function<bool(const Registration&)>;

bool validate_registration(const Registration& registration, std::span<const Rule> rules)
{
    for (const Rule& rule : rules) {
        if (!rule(registration)) {
            return false;
        }
    }
    return true;
}
\end{lstlisting}

组装规则：

\begin{lstlisting}[language=C++,caption={组装带状态规则}]
std::vector<Rule> rules;
rules.push_back([](const Registration& item) {
    return is_valid_username(item.username);
});
rules.push_back([](const Registration& item) {
    return item.age >= 18;
});

std::string allowed_domain = "@example.com";
rules.push_back([allowed_domain](const Registration& item) {
    return item.email.ends_with(allowed_domain);
});
\end{lstlisting}

这里 \code{allowed_domain} 按值捕获，因为规则会保存在 \code{rules} 里，可能比局部变量活得更久。这个小案例把 lambda、捕获、\code{std::function} 和生命周期放到同一个场景里。

\begin{exercisebox}
给验证流水线增加错误消息，不要只返回 \code{bool}。设计 \code{RuleResult}，让每条规则失败时返回原因。要求规则仍然可组合，验证函数返回第一条错误或全部错误列表。
\end{exercisebox}
